\subsection{Программная реализация системы}
\label{sec:development:implementation}

\subsubsection{Архитектура системы}

Система состоит из нескольких сервисов, каждый из которых запущен в отдельном контейнере на одном сервере. Архитектура развёртывания приведена на рисунке~\ref{fig:development:architecture}.

\begin{figure}[H]
	\centering
	\includegraphics[width=\linewidth]{kufar/architecture}
	\caption{Архитектура развёртывания системы}
	\label{fig:development:architecture}
\end{figure}

Основную работу выполняют три сервиса. СУБД \textit{PostgreSQL} хранит все данные. Серверный сервис владеет базой данных: только этот сервис выполняет SQL-запросы, применяет миграции и запускает периодические задачи. Бот отвечает только за диалог с пользователем в \textit{Telegram}. Собственного доступа к базе данных у бота нет: все данные для показа и изменения бот получает от серверного сервиса по \textit{HTTP}, прикладывая к каждому запросу общий секретный ключ.

Такое разделение выбрано ради данных. Когда в базу пишет только один сервис, все правила записи собраны в одном месте. Лимит отслеживаемых товаров, например, проверяется в серверном сервисе, и обойти его нельзя ни через бот, ни через импорт файла, ни прямым запросом к интерфейсу. Схемой базы данных управляет тот же сервис, поэтому несогласованных изменений структуры не бывает.

Уведомления серверный сервис отправляет в \textit{Telegram} сам, не через бота. Для отправки сообщений достаточно токена бота, а получать входящие сообщения продолжает сервис бота. Поэтому остановка бота, например при обновлении, не прерывает доставку уведомлений.

Вспомогательные сервисы обеспечивают надёжность. Сервис резервного копирования раз в сутки выгружает дамп базы. Служба мониторинга следит за событиями контейнеров и сообщает администратору о запуске, падении и потере работоспособности любого из них. События \textit{Docker} служба читает через отдельный посредник, который открывает к ним доступ только на чтение: прямой доступ к управляющему сокету \textit{Docker} равносилен правам администратора сервера. Почтовый сервер-ловушка принимает все письма системы и показывает их в веб-интерфейсе, что упрощает проверку почтовых функций без настоящего почтового ящика.

Внутри серверного сервиса код разделён на слои. Модули обработки и их связи с сущностями данных показаны на рисунке~\ref{fig:development:classes:logic}.

Сущности данных описывают содержимое таблиц базы данных:
\begin{itemize}
	\item \textit{User}~-- пользователь: номер учётной записи \textit{Telegram}, язык интерфейса, признаки активности и уведомлений, адрес электронной почты и признак его подтверждения;
	\item \textit{TrackedProduct}~-- отслеживаемый товар: название, категория, границы цены, списки слов поиска и исключений, фильтры площадки в поле \textit{kufar\_filters} и порог выгодности \textit{min\_deal\_score};
	\item \textit{Listing}~-- объявление, полученное с площадки: площадка, внешний идентификатор, заголовок, описание, цена, ссылка, время публикации, время первого и последнего появления в выдаче;
	\item \textit{Match}~-- совпадение объявления с отслеживаемым товаром: показатель выгоды, подробности отбора, время доставки \textit{notified\_at}, число попыток \textit{notify\_attempts} и текст ошибки;
	\item \textit{SentNotification}~-- запись журнала отправленных уведомлений с ценой объявления на момент отправки;
	\item \textit{KufarAttribute} и \textit{KufarAttributeValue}~-- справочник характеристик категорий площадки и их допустимых значений.
\end{itemize}

Доступ к площадке построен на наследовании. Абстрактный класс \textit{BaseScraper} хранит сетевую сессию, удерживает частоту запросов методом \textit{\_rate\_limit} и выполняет запрос методом \textit{\_get}, а поиск и загрузку отдельного объявления объявляет, но не реализует. Класс \textit{KufarScraper} реализует их для \textit{Kufar.by}: метод \textit{\_build\_params} собирает параметры запроса по стратегии категории, методы \textit{\_parse\_search\_response} и \textit{\_parse\_ad} превращают ответ площадки в структуры \textit{ScrapedListing}. Добавление второй площадки сводится к новому наследнику \textit{BaseScraper}, потому что остальной код работает со структурой \textit{ScrapedListing} и об устройстве площадки ничего не знает.

Обработка данных вынесена в отдельные модули и классы, показанные на рисунке~\ref{fig:development:classes:logic}.

\begin{figure}[H]
	\centering
	\includegraphics[width=\linewidth]{kufar/classes_logic}
	\caption{Диаграмма классов модулей обработки}
	\label{fig:development:classes:logic}
\end{figure}

Назначение модулей:
\begin{itemize}
	\item \textit{engine}~-- функция \textit{persist\_listing} добавляет или обновляет объявление и возвращает прежнюю цену, функция \textit{match\_listings\_for\_product} отбирает объявления по критериям товара, создаёт совпадения и возвращает список событий \textit{MatchEvent};
	\item \textit{sender}~-- функция \textit{deliver\_pending} разбирает очередь недоставленных совпадений, функция \textit{send\_deal\_notification} отправляет одно уведомление, функция \textit{\_min\_sent\_price} проверяет журнал на повтор, функция \textit{give\_up\_on\_stale} закрывает совпадения старше трёх суток;
	\item \textit{limits}~-- функция \textit{check\_can\_add} проверяет лимит отслеживаемых товаров и возвращает структуру \textit{LimitStatus}, функции \textit{is\_admin} и \textit{limit\_for} определяют потолок для конкретного пользователя;
	\item \textit{jobs}~-- периодические задачи планировщика: опрос площадки \textit{scrape\_and\_match\_job}, доставка уведомлений \textit{deliver\_notifications\_job}, почтовые копии \textit{deliver\_email\_copies\_job}, очистка устаревших данных \textit{cleanup\_old\_data\_job} и обновление справочника \textit{refresh\_kufar\_schema\_job}.
\end{itemize}

В структуре \textit{CategoryStrategy} собраны правила отбора для каждой категории: код категории площадки, обязательные фильтры, порядок характеристик и их числовые диапазоны. Результаты работы модули передают друг другу структурами: \textit{ScrapedListing}~-- разобранное объявление площадки, \textit{MatchEvent}~-- найденное совпадение вместе с прежней ценой и типом события, \textit{LimitStatus}~-- результат проверки лимита. На диаграммах не показаны части системы, которые не работают с данными напрямую: обработчики бота, маршруты программного интерфейса, словари локализации и формирование книг \textit{Excel}.

Слои обращаются друг к другу сверху вниз. Маршрут программного интерфейса проверяет входные данные и вызывает сервис бизнес-логики. Сервис выполняет SQL-запросы к базе данных и возвращает результат. Модуль \textit{jobs} обращается к тем же модулям \textit{engine}, \textit{sender} и \textit{limits}, что и маршруты, поэтому одно и то же правило, например проверка повтора уведомления, реализовано один раз. Класс \textit{KufarScraper} о базе данных ничего не знает и возвращает объявления структурами \textit{ScrapedListing}, а сохраняет их функция \textit{persist\_listing}.

\subsubsection{Цикл опроса площадки и обработка объявлений}

Цикл опроса запускается планировщиком каждые 3 минуты. Одновременно может выполняться только один цикл: если предыдущий ещё не завершился, следующий запуск пропускается. Цикл выполняет следующие шаги:
\begin{enumerate}
	\item Проверить паузу. Если после отказа площадки ещё остались пропускаемые циклы, уменьшить их счётчик и завершить цикл.
	\item Выбрать из базы данных все активные отслеживаемые товары.
	\item Для каждого товара по очереди:
	\begin{itemize}[itemindent=2.2ex,leftmargin=2\parindent,label=--]
		\item пропустить товар, если у него не задана пороговая цена;
		\item по стратегии категории сформировать запрос: код категории площадки, тип сделки, обязательные фильтры, фильтры пользователя, название и ограничение цены сверху, равное порогу;
		\item дождаться разрешения ограничителя частоты и отправить запрос;
		\item если площадка ответила кодом 403 или 429, прервать цикл целиком;
		\item передать полученные объявления модулю отбора.
	\end{itemize}
	\item Зафиксировать транзакцию.
	\item Если цикл занял больше 80\,\% интервала, записать предупреждение в журнал работы.
\end{enumerate}

Ограничитель частоты работает в два слоя. Перед каждым запросом выдерживается пауза около 2 секунд со случайной добавкой от 0,5 до 1,5 секунды, чтобы запросы не шли с одинаковым интервалом. Кроме того, ограничитель помнит время запросов за последнюю минуту и, если их уже 30, ждёт, пока самый старый выйдет из этого окна. При отказе площадки пауза до следующего опроса удваивается с каждым отказом подряд: 1, 2, 4, 8, 16 циклов. Администратор получает одно оповещение, а повторные оповещения того же вида в течение 30 минут подавляются.

Модуль отбора обрабатывает каждое полученное объявление так:
\begin{enumerate}
	\item Найти объявление в базе данных по паре «площадка~-- внешний идентификатор».
	\item Если объявление найдено, запомнить его прежнюю цену и обновить запись новыми данными, иначе добавить новую запись.
	\item Проверить релевантность. Объявление отклоняется, если выполняется хотя бы одно условие:
	\begin{itemize}[itemindent=2.2ex,leftmargin=2\parindent,label=--]
		\item цена отсутствует, не ниже порога или ниже минимальной правдоподобной цены категории;
		\item заголовок содержит стоп-слово категории (для электроники это «чехол», «защитное стекло» и подобные);
		\item заголовок или описание содержат слово-исключение пользователя;
		\item в заголовке или описании нет какого-либо слова из названия товара (в категориях с фильтрами площадки проверка названия не выполняется);
		\item нет какого-либо обязательного слова.
	\end{itemize}
	\item Если совпадение для этой пары «товар~-- объявление» уже есть, обновить его только при снижении цены.
	\item Иначе определить тип события: «новое», если объявление опубликовано после добавления товара, «снижение цены», если цена упала с прошлого опроса. Если ни то ни другое, объявление пропускается.
	\item Создать совпадение с показателем выгоды и пустым временем доставки.
\end{enumerate}

Проверка названия допускает расхождения в написании. Слово из названия ищется в три этапа: сначала как подстрока текста объявления, затем по словарю известных вариантов записи (латиница и кириллица, распространённые сокращения), затем нечётким сравнением с каждым словом текста. На последнем этапе слово длиной от трёх символов считается найденным, если мера сходства на основе редакционного расстояния составляет не менее 85\,\%.

Отдельно работает мгновенный поиск. Сразу после добавления товара система выполняет тот же запрос к площадке и те же проверки релевантности, сортирует прошедшие объявления по цене и показывает пользователю пять самых дешёвых. В базу данных эти объявления не записываются: мгновенный поиск показывает состояние рынка, а отслеживание начинается с момента добавления товара.

\subsubsection{Механизм доставки уведомлений}

Уведомления доставляются по схеме исходящей очереди (\textit{outbox}). Цикл опроса только записывает совпадения в базу данных, а отправляет их отдельная задача, которую планировщик запускает каждую минуту. Очередью служит сама таблица совпадений: в очереди находятся записи с пустым временем доставки. Отдельной таблицы или брокера сообщений для этого не нужно: запись совпадения и постановка в очередь выполняются одной операцией в одной транзакции.

Задача доставки выполняет следующие шаги:
\begin{enumerate}
	\item Выбрать до 50 совпадений в порядке создания, для которых одновременно:
	\begin{itemize}[itemindent=2.2ex,leftmargin=2\parindent,label=--]
		\item время доставки пусто;
		\item неудачных попыток меньше пяти;
		\item совпадение создано не раньше трёх суток назад;
		\item товар активен, пользователь активен и не отключил уведомления.
	\end{itemize}
	\item Сгруппировать совпадения по товарам.
	\item Для каждого товара посчитать, сколько уведомлений ещё можно отправить: не больше пяти за последний час по этому товару и не больше 30 за текущие сутки пользователю. Если лимит исчерпан, оставить совпадения в очереди.
	\item Отсортировать совпадения товара по цене, чтобы при ограниченном лимите первыми ушли самые выгодные.
	\item Для каждого совпадения найти в журнале минимальную цену, по которой это объявление уже отправлялось пользователю:
	\begin{itemize}[itemindent=2.2ex,leftmargin=2\parindent,label=--]
		\item если такой отправки не было, отправить уведомление;
		\item если новая цена ниже минимальной, отправить уведомление с пометкой о снижении цены;
		\item иначе закрыть совпадение как дубликат.
	\end{itemize}
	\item После успешной отправки заполнить время доставки и добавить запись в журнал уведомлений.
	\item Зафиксировать транзакцию.
\end{enumerate}

Ошибки отправки обрабатываются по-разному. Если \textit{Telegram} просит снизить частоту запросов, доставка прекращается до следующего запуска: это ограничение действует на бота целиком, и продолжать отправку другим пользователям бессмысленно. Пользователь, заблокировавший бота, помечается неактивным, и его совпадения перестают выбираться из очереди. При любой другой ошибке у совпадения увеличивается счётчик попыток и сохраняется текст ошибки, а само совпадение остаётся в очереди до следующего запуска. После пятой неудачной попытки совпадение из очереди исключается, а число таких совпадений записывается в журнал работы.

Почтовые копии ведут отдельную очередь в той же таблице совпадений, со своим временем отправки и счётчиком попыток. Копия отправляется только для совпадения, уже доставленного в \textit{Telegram}, поэтому все проверки лимитов и повторов действуют и для почты. Раз в 5 минут задача собирает все новые доставленные совпадения пользователя, у которого подтверждён адрес и включены копии, и отправляет одно сводное письмо со списком до 20 объявлений. Сбой почтового сервера не задерживает доставку в \textit{Telegram}, потому что две очереди обрабатываются разными задачами.

\subsubsection{Лимиты отслеживания и права администратора}

Каждый отслеживаемый товар добавляет один запрос к площадке за цикл опроса, поэтому число товаров ограничено двумя пределами:
\begin{itemize}
	\item личный лимит: 3 товара для пользователя и 10 для администратора;
	\item общий лимит: не более 25 товаров на всю систему.
\end{itemize}

Личный лимит защищает от ситуации, когда один пользователь занимает всю систему. Общий лимит защищает площадку и саму систему: при интервале опроса 180 секунд и примерно 3 секундах на запрос с учётом паузы за цикл можно обработать около 60 товаров. Предел в 25 товаров оставляет запас, чтобы цикл гарантированно завершался до начала следующего. При запуске серверный сервис сравнивает общий лимит с этой расчётной ёмкостью и предупреждает в журнале работы, если лимит задан больше, чем цикл успевает обработать.

Проверка лимитов выполняется дважды: перед началом диалога добавления товара и в момент сохранения. Вторая проверка нужна потому, что за время диалога другой пользователь мог занять последнее свободное место. Логика проверки такова:
\begin{enumerate}
	\item Посчитать товары пользователя и товары всех пользователей.
	\item Если товаров пользователя не меньше его личного лимита, отказать с пояснением «достигнут личный лимит».
	\item Если товаров в системе не меньше 25, отказать с пояснением «достигнут общий лимит системы».
	\item Иначе разрешить добавление.
\end{enumerate}

Все пути добавления товара проходят через одну и ту же операцию сервиса товаров: пошаговый диалог, импорт из ссылки и импорт файла \textit{CSV}. При импорте файла товары добавляются по строкам, пока лимит не исчерпан. В отчёте об импорте пользователь видит число добавленных товаров, ошибки по номерам строк и сообщение о том, какой лимит остановил импорт.

Права администратора определяются списком идентификаторов \textit{Telegram} в конфигурации системы, отдельной таблицы ролей нет. Администраторов в системе единицы, и назначает их владелец сервера, изменяя конфигурацию. Администратор отличается от пользователя тремя возможностями: личным лимитом в 10 товаров, отчётом о состоянии системы и служебными оповещениями о блокировке со стороны площадки и о событиях контейнеров. Остальные функции у него те же, что у пользователя, и доступа к чужим товарам администратор не получает: любая операция с товаром выполняется только в пределах товаров владельца.
